\SourceLeaf{142}
\[
\mathrm{VI.}\quad 0=\gamma F-\gamma(1-x)F''-(\gamma-\alpha-\beta-1)xF'''
\]
Hinc atque ex III, eliminando $F''$
\[
\mathrm{VII.}\quad 0=\gamma F-(\gamma-\alpha-1-\beta x)F'''-(\alpha+1)(1-x)F''''
\]
Porro ex IV atque V, eliminando $F'''''$
\[
\mathrm{VIII.}\quad 0=(\gamma+1)F'''-(\gamma+1)F''''+(\beta+1)xF''''''
\]
Hinc atque ex VII, eliminando $F''''$,
\[
\begin{aligned}&\mathrm{IX.}\quad 0=\gamma(\gamma+1)F\\&{}-(\gamma+1)(\gamma-(\alpha+\beta+1)x)F'''\\&{}-(\alpha+1)(\beta+1)x(1-x)F''''''\end{aligned}
\]
\SixthArticle{11}
Si omnes relationes inter ternas functiones $F(\alpha,\beta,\gamma)$, $F(\alpha+\lambda,\beta+\mu,\gamma+\nu)$, $F(\alpha+\lambda',\beta+\mu',\gamma+\nu')$, in quibus $\lambda,\mu,\nu,\lambda',\mu',\nu'$ vel $=0$ vel $=+1$ vel $=-1$, exhaurire vellemus, formularum multitudo usque ad 325 ascenderet. Haud inutilis foret talis collectio, saltem simpliciorum ex his formulis: hoc vero loco sufficiat, paucas tantummodo apposuisse, quas vel ex formulis art.\ 7, vel si magis placet, simili modo ut duae priores ex illis in art.\ 8 erutae sunt, quivis nullo negotio sibi demonstrare poterit.
\[
\begin{aligned}&[16]\quad F(\alpha,\beta,\gamma)-F(\alpha,\beta,\gamma-1)\\&{}=-\frac{\alpha\beta x}{\gamma(\gamma-1)}F(\alpha+1,\beta+1,\gamma+1)
\end{aligned}
\]
\[
\begin{aligned}&[17]\quad F(\alpha,\beta+1,\gamma)-F(\alpha,\beta,\gamma)\\&{}=\frac{\alpha x}{\gamma}F(\alpha+1,\beta+1,\gamma+1)
\end{aligned}
\]
\[
\begin{aligned}&[18]\quad F(\alpha+1,\beta,\gamma)-F(\alpha,\beta,\gamma)\\&{}=\frac{\beta x}{\gamma}F(\alpha+1,\beta+1,\gamma+1)
\end{aligned}
\]
\[
\begin{aligned}&[19]\quad F(\alpha,\beta+1,\gamma+1)-F(\alpha,\beta,\gamma)\\&{}=\frac{\alpha(\gamma-\beta)x}{\gamma(\gamma+1)}F(\alpha+1,\beta+1,\gamma+2)
\end{aligned}
\]
\[
\begin{aligned}&[20]\quad F(\alpha+1,\beta,\gamma+1)-F(\alpha,\beta,\gamma)\\&{}=\frac{\beta(\gamma-\alpha)x}{\gamma(\gamma+1)}F(\alpha+1,\beta+1,\gamma+2)
\end{aligned}
\]
\[
\begin{aligned}&[21]\quad F(\alpha-1,\beta+1,\gamma)-F(\alpha,\beta,\gamma)\\&{}=\frac{(\alpha-\beta-1)x}{\gamma}F(\alpha,\beta+1,\gamma+1)
\end{aligned}
\]
\[
\begin{aligned}&[22]\quad F(\alpha+1,\beta-1,\gamma)-F(\alpha,\beta,\gamma)\\&{}=\frac{(\beta-\alpha-1)x}{\gamma}F(\alpha+1,\beta,\gamma+1)
\end{aligned}
\]
\[
\begin{aligned}&[23]\quad F(\alpha-1,\beta+1,\gamma)-F(\alpha+1,\beta-1,\gamma)\\&{}=\frac{(\alpha-\beta)x}{\gamma}F(\alpha+1,\beta+1,\gamma+1)
\end{aligned}
\]
\begin{center}\rule{.24\linewidth}{.3pt}\end{center}
\SourceLeaf{143}
\par\Needspace{9\baselineskip}\pdfbookmark[1]{Fractiones continuae}{sixth-section-2}\begin{center}\textbf{\textit{SECTIO SECUNDA.}}\\[6pt]\textbf{\textit{Fractiones continuae.}}\end{center}
\SixthArticle{12}
Designando
\[
\frac{F(\alpha,\beta+1,\gamma+1,x)}{F(\alpha,\beta,\gamma,x)}\quad\hbox{per}\quad G(\alpha,\beta,\gamma,x)
\]
fit
\[
\frac{F(\alpha+1,\beta,\gamma+1,x)}{F(\alpha,\beta,\gamma,x)}=\frac{F(\beta,\alpha+1,\gamma+1,x)}{F(\beta,\alpha,\gamma,x)}=G(\beta,\alpha,\gamma,x)
\]
et proin, dividendo aequationem 19 per $F(\alpha,\beta+1,\gamma+1,x)$,
\[
1-\frac{1}{G(\alpha,\beta,\gamma,x)}=\frac{\alpha(\gamma-\beta)}{\gamma(\gamma+1)}x\,G(\beta+1,\alpha,\gamma+1,x)
\]
sive
\[
[24]\qquad G(\alpha,\beta,\gamma,x)=\cfrac{1}{1-\frac{\alpha(\gamma-\beta)}{\gamma(\gamma+1)}x\,G(\beta+1,\alpha,\gamma+1,x)}
\]
et quum perinde fiat
\[
G(\beta+1,\alpha,\gamma+1,x)=\cfrac{1}{1-\frac{(\beta+1)(\gamma+1-\alpha)}{(\gamma+1)(\gamma+2)}x\,G(\alpha+1,\beta+1,\gamma+2,x)}
\]
etc., resultabit pro $G(\alpha,\beta,\gamma,x)$ fractio continua
\[
[25]\qquad\frac{F(\alpha,\beta+1,\gamma+1,x)}{F(\alpha,\beta,\gamma,x)}=\cfrac{1}{1-\cfrac{ax}{1-\cfrac{bx}{1-\cfrac{cx}{1-\cfrac{dx}{1-\hbox{etc.}}}}}}
\]
ubi
\[
\begin{array}{ll}
a=\frac{\alpha(\gamma-\beta)}{\gamma(\gamma+1)}&b=\frac{(\beta+1)(\gamma+1-\alpha)}{(\gamma+1)(\gamma+2)}\\[6pt]
c=\frac{(\alpha+1)(\gamma+1-\beta)}{(\gamma+2)(\gamma+3)}&d=\frac{(\beta+2)(\gamma+2-\alpha)}{(\gamma+3)(\gamma+4)}\\[6pt]
e=\frac{(\alpha+2)(\gamma+2-\beta)}{(\gamma+4)(\gamma+5)}&f=\frac{(\beta+3)(\gamma+3-\alpha)}{(\gamma+5)(\gamma+6)}
\end{array}
\]
etc., cuius lex progressionis obvia est.

Porro ex aequationibus 17, 18, 21, 22 sequitur
\SourceLeaf{144}
\[
\begin{array}{rl}
[26]&\frac{F(\alpha,\beta+1,\gamma,x)}{F(\alpha,\beta,\gamma,x)}=\cfrac{1}{1-\frac{\alpha x}{\gamma}G(\beta+1,\alpha,\gamma,x)}\\[12pt]
[27]&\frac{F(\alpha+1,\beta,\gamma,x)}{F(\alpha,\beta,\gamma,x)}=\cfrac{1}{1-\frac{\beta x}{\gamma}G(\alpha+1,\beta,\gamma,x)}\\[12pt]
[28]&\frac{F(\alpha-1,\beta+1,\gamma,x)}{F(\alpha,\beta,\gamma,x)}=\cfrac{1}{1-\frac{(\alpha-\beta-1)x}{\gamma}G(\beta+1,\alpha-1,\gamma,x)}\\[12pt]
[29]&\frac{F(\alpha+1,\beta-1,\gamma,x)}{F(\alpha,\beta,\gamma,x)}=\cfrac{1}{1-\frac{(\beta-\alpha-1)x}{\gamma}G(\alpha+1,\beta-1,\gamma,x)}
\end{array}
\]
unde, substitutis pro functione $G$ eius valoribus in fractionibus continuis, totidem fractiones continuae novae prodeunt.

Ceterum sponte patet, fractionem continuam in formula 25 abrumpi, si e numeris $\alpha,\beta,\gamma-\alpha,\gamma-\beta$ aliquis fuerit integer negativus, alioquin vero in infinitum excurrere.
\SixthArticle{13}
Fractiones continuae in art.\ praec.\ erutae maximi sunt momenti, asserique potest, vix ullas fractiones continuas secundum legem obviam progredientes ab analystis hactenus erutas esse, quae sub nostris tamquam casus speciales non sint contentae. Imprimis memorabilis est casus is, ubi in formula 25 statuitur $\beta=0$, unde $F(\alpha,\beta,\gamma,x)=1$, adeoque, scribendo $\gamma-1$ pro $\gamma$
\[
\begin{aligned}
{}[30]\quad F(\alpha,1,\gamma)&=1+\frac{\alpha}{\gamma}x+\frac{\alpha(\alpha+1)}{\gamma(\gamma+1)}xx\\&{}+\frac{\alpha(\alpha+1)(\alpha+2)}{\gamma(\gamma+1)(\gamma+2)}x^3+\hbox{etc.}\\[9pt]
&=\cfrac{1}{1-\cfrac{ax}{1-\cfrac{bx}{1-\cfrac{cx}{1-\cfrac{dx}{1-\hbox{etc.}}}}}}
\end{aligned}
\]
ubi
\[
\begin{array}{ll}
a=\frac{\alpha}{\gamma}&b=\frac{\gamma-\alpha}{\gamma(\gamma+1)}\\[6pt]
c=\frac{(\alpha+1)\gamma}{(\gamma+1)(\gamma+2)}&d=\frac{2(\gamma+1-\alpha)}{(\gamma+2)(\gamma+3)}\\[6pt]
e=\frac{(\alpha+2)(\gamma+1)}{(\gamma+3)(\gamma+4)}&f=\frac{3(\gamma+2-\alpha)}{(\gamma+4)(\gamma+5)}
\end{array}
\]
\begin{center}etc.\end{center}
\SourceLeaf{145}
\SixthArticle{14}

Operae pretium erit, quosdam casus speciales huc adscripsisse. Ex formula I art. 5 sequitur, statuendo $t=1$, $\beta=1$
\[
[31]\qquad
(1+u)^n=
\cfrac{1}{1-\cfrac{nu}{1+\cfrac{\frac{n+1}{2}u}{1-\cfrac{\frac{n-1}{2\cdot3}u}{1+\cfrac{\frac{2(n+2)}{3\cdot4}u}{1-\cfrac{\frac{2(n-2)}{4\cdot5}u}{1+\mathrm{etc.}}}}}}}
\]

E formulis VI, VII art. 5 sequitur
\[
[32]\qquad
\log(1+t)=
\cfrac{t}{1+\cfrac{\frac12t}{1+\cfrac{\frac16t}{1+\cfrac{\frac26t}{1+\cfrac{\frac2{10}t}{1+\cfrac{\frac3{10}t}{1+\cfrac{\frac3{14}t}{1+\mathrm{etc.}}}}}}}}
\]
\[
[33]\qquad
\log\cfrac{1+t}{1-t}=
\cfrac{2t}{1-\cfrac{\frac13tt}{1-\cfrac{\frac{2\cdot2}{3\cdot5}tt}{1-\cfrac{\frac{3\cdot3}{5\cdot7}tt}{1-\cfrac{\frac{4\cdot4}{7\cdot9}tt}{1-\mathrm{etc.}}}}}}
\]
Mutando hic signa $-$ in $+$ prodit fractio continua pro arc. tang $t$.
\par
Porro habemus
\[
[34]\qquad
e^t=
\cfrac{1}{1-\cfrac{t}{1+\cfrac{\frac13t}{1-\cfrac{\frac16t}{1+\cfrac{\frac16t}{1-\cfrac{\frac1{10}t}{1+\cfrac{\frac1{10}t}{1-\mathrm{etc.}}}}}}}}
\]
\SourceLeaf{146}
\[
[35]\qquad
t=
\cfrac{\sin t\cos t}{1-\cfrac{\frac{1\cdot2}{1\cdot3}\sin t^2}{1-\cfrac{\frac{1\cdot2}{3\cdot5}\sin t^2}{1-\cfrac{\frac{3\cdot4}{5\cdot7}\sin t^2}{1-\cfrac{\frac{3\cdot4}{7\cdot9}\sin t^2}{1-\cfrac{\frac{5\cdot6}{9\cdot11}\sin t^2}{1-\mathrm{etc.}}}}}}}
\]

Statuendo $\alpha=3$, $\gamma=\frac52$, e formula 30 sponte sequitur fractio continua in \textit{Theoria motus corporum coelestium} art. 90 proposita. Ibidem duae aliae fractiones continuae prolatae sunt, quarum evolutionem hacce occasione supplere visum est. Statuendo
\[
Q=1-\cfrac{\frac{5\cdot8}{7\cdot9}x}{1-\cfrac{\frac{1\cdot4}{9\cdot11}x}{1-\frac{7\cdot10}{11\cdot13}x\ \mathrm{etc.}}}
\]
fit l. c. $\displaystyle x-\xi=\cfrac{x}{1+\frac{2x}{35Q}}=\frac{xQ}{Q+\frac2{35}x}$, adeoque
\[
\xi=\frac{\frac2{35}xx}{Q+\frac2{35}x}
\]
quae est formula prior: posterior sequenti modo eruitur. Statuendo
\[
R=1-\cfrac{\frac{1\cdot4}{7\cdot9}x}{1-\cfrac{\frac{5\cdot8}{9\cdot11}x}{1-\cfrac{\frac{3\cdot6}{11\cdot13}x}{1-\frac{7\cdot10}{13\cdot15}x\ \mathrm{etc.}}}}
\]
erit per formulam 25
\[
\frac1R=G(\frac12,\frac32,\frac72,x),
\qquad\mbox{atque}\qquad
\frac1Q=G(\frac52,-\frac12,\frac72,x)
\]
Hinc
\SourceLeaf{147}
\[
\begin{array}{rcl}
RF(\frac12,\frac52,\frac92,x)&=&F(\frac12,\frac32,\frac72,x)\\[4pt]
QF(\frac52,\frac12,\frac92,x)&=&F(\frac52,-\frac12,\frac72,x)
\end{array}
\]
sive permutando elementum primum cum secundo
\[
QF(\frac12,\frac52,\frac92,x)=F(-\frac12,\frac52,\frac72,x)
\]
\par
Sed per aequationem 21 habemus
\[
F(-\frac12,\frac52,\frac72,x)-F(\frac12,\frac32,\frac72,x)
=-\frac47xF(\frac12,\frac52,\frac92,x)
\]
unde fit $Q=R-\frac47x$, quo valore in formula supra data substituto prodit
\[
\xi=\frac{\frac2{35}xx}{R-\frac{18}{35}x}
\]
quae est formula posterior.
\par
Statuendo in formula 30, $\alpha=\frac mn$, $x=-\gamma nt$, fit pro valore infinite magno ipsius $\gamma$
\[
\begin{aligned}
{}[36]\qquad F(\frac mn,1,\gamma,-\gamma nt)
&=1-mt+m(m+n)tt\\&{}-m(m+n)(m+2n)t^3+\mathrm{etc.}\\[8pt]
&=\displaystyle\cfrac{1}{1+\cfrac{mt}{1+\cfrac{nt}{1+\cfrac{(m+n)t}{1+\cfrac{2nt}{1+\cfrac{(m+2n)t}{1+3nt\ \mathrm{etc.}}}}}}}
\end{aligned}
\]
\begin{center}\rule{0.27\textwidth}{0.4pt}\end{center}
\par\Needspace{9\baselineskip}\pdfbookmark[1]{De summa seriei}{sixth-section-3}\begin{center}
\textit{SECTIO TERTIA.}\\[6pt]
\textit{De summa seriei nostrae statuendo elementum quartum $=1$, ubi simul quaedam aliae functiones}\\
\textit{transscendentes discutiuntur.}\\[6pt]
\end{center}
\SixthArticle{15}

Quoties elementa $\alpha,\beta,\gamma$ omnia sunt quantitates positivae, omnes co\"efficientes potestatum elementi quarti $x$ positivi evadunt: quoties vero ex illis elementis unum alterumve negativum est, saltem inde ab aliqua potestate $x^m$ omnes co\"efficientes eodem signo affecti erunt, si modo $m$ accipitur maior quam valor absolutus elementi negativi maximi. Porro hinc sponte patet, seriei summam pro
\SourceLeaf{148}
$x=1$ finitam esse non posse, nisi co\"efficientes saltem post certum terminum in infinitum decrescant, vel, ut secundum morem analystarum loquamur, nisi co\"efficiens termini $x^{\infty}$ sit $=0$. Ostendemus autem, et quidem, in gratiam eorum, qui methodis rigorosis antiquorum geometrarum favent, omni rigore,

\textit{primo}, co\"efficientes (siquidem series non abrumpatur), in infinitum crescere, quoties fuerit $\alpha+\beta-\gamma-1$ quantitas positiva.

\textit{secundo}, co\"efficientes versus limitem finitum continuo convergere, quoties fuerit $\alpha+\beta-\gamma-1=0$.

\textit{tertio}, co\"efficientes in infinitum decrescere, quoties fuerit $\alpha+\beta-\gamma-1$ quantitas negativa.

\textit{quarto}, summam seriei nostrae pro $x=1$, non obstante convergentia in casu tertio, infinitam esse, quoties fuerit $\alpha+\beta-\gamma$ quantitas positiva vel $=0$.

\textit{quinto}, summam vero \textit{finitam} esse, quoties $\alpha+\beta-\gamma$ fuerit quantitas negativa.

\SixthArticle{16}
Hanc disquisitionem generalius adaptabimus seriei infinitae $M,\allowbreak M',\allowbreak M'',\allowbreak M'''$ etc. ita formatae, ut quotientes $\frac{M'}{M},\frac{M''}{M'},\frac{M'''}{M''}$ etc. resp. sint valores fractionis
\[
\frac{t^{\lambda}+At^{\lambda-1}+Bt^{\lambda-2}+Ct^{\lambda-3}+\mathrm{etc.}}
{t^{\lambda}+at^{\lambda-1}+bt^{\lambda-2}+ct^{\lambda-3}+\mathrm{etc.}}
\]
pro $t=m$, $t=m+1$, $t=m+2$ etc. Brevitatis caussa huius fractionis numeratorem per $P$, denominatorem per $p$ denotabimus: praeterea supponemus, $P,p$ non esse identicas, sive differentias $A-a,B-b,C-c$ etc. non omnes simul evanescere.

I. Quoties e differentiis $A-a,B-b,C-c$ etc. prima quae non evanescit est positiva, assignari poterit limes aliquis $l$, quem simulac egressus est valor ipsius $t$, valores functionum $P$ et $p$ certo semper evadent positivi, atque $P>p$. Manifestum est, hoc evenire, si pro $l$ accipiatur radix maxima realis aequationis $p(P-p)=0$; si vero haec aequatio nullas omnino radices reales habeat, proprietatem illam pro omnibus valoribus ipsius $t$ locum habere. Quapropter in serie $\frac{M'}{M},\frac{M''}{M'},\frac{M'''}{M''}$ etc. saltem post certum intervallum (si non ab initio) omnes termini erunt positivi atque maiores unitate; quodsi itaque nullus neque $=0$ neque infinite magnus evadit, perspicuum est,
\SourceLeaf{149}
\textit{seriem $M,M',M'',M'''$ etc. si non ab initio tamen post certum intervallum omnes suos terminos eodem signo affectos continuoque crescentes habituram esse.}

Eadem ratione, si e differentiis $A-a,B-b,C-c$ etc. prima quae non evanescit est negativa, series $M,M',M'',M'''$ etc. si non ab initio tamen post certum intervallum omnes suos terminos eodem signo affectos continuoque decrescentes habebit.

II. Si iam co\"efficientes $A,a$ sunt inaequales, termini seriei $M,M',M'',M'''$ etc. ultra omnes limites sive in infinitum vel crescent vel decrescent, prout differentia $A-a$ est positiva vel negativa: hoc ita demonstramus. Si $A-a$ est quantitas positiva, accipiatur numerus integer $h$ ita, ut fiat $h(A-a)>1$, statuaturque $\frac{M^h}{m}=N$, $\frac{M^{\prime h}}{m+1}=N'$, $\frac{M^{\prime\prime h}}{m+2}=N''$, $\frac{M^{\prime\prime\prime h}}{m+3}=N'''$ etc., nec non $tP^h=Q$, $(t+1)p^h=q$. Tunc patet, $\frac{N'}{N},\frac{N''}{N'},\frac{N'''}{N''}$ etc. esse valores fractionis $\frac{Q}{q}$ ponendo $t=m$, $t=m+1$, $t=m+2$ etc., ipsas $Q,q$ vero esse functiones algebraicas formae huius
\[
\begin{array}{l}
Q=t^{\lambda h+1}+hAt^{\lambda h}+\mathrm{etc.}\\
q=t^{\lambda h+1}+(ha+1)t^{\lambda h}+\mathrm{etc.}
\end{array}
\]
Quare quum per hyp. differentia $hA-(ha+1)$ sit quantitas positiva, termini seriei $N,N',N'',N'''$ etc. si non ab initio tamen post certum intervallum continuo crescent (per \textsc{i}); hinc termini seriei $mN,(m+1)N',(m+2)N''$, $(m+3)N'''$ etc. necessario ultra omnes limites crescent, et proin etiam termini seriei $M,M',M'',M'''$ etc., quippe quorum potestates exponente $h$ illis sunt aequales. Q.~E.~P.

Si $A-a$ est quantitas negativa, accipere oportet integrum $h$ ita, ut $h(a-A)$ fiat maior quam $1$, unde per ratiocinia similia termini seriei
\[
mM^h,\quad (m+1)M^{\prime h},\quad (m+2)M^{\prime\prime h},\quad (m+3)M^{\prime\prime\prime h}\quad\mathrm{etc.}
\]
post certum intervallum continuo decrescent. Quamobrem termini seriei $M^h$, $M^{\prime h},M^{\prime\prime h}$ etc. adeoque etiam termini huius $M,M',M'',M'''$ etc. necessario in infinitum decrescent. Q.~E.~S.

III. Si vero co\"efficientes primi $A,a$ sunt aequales, termini seriei $M,\allowbreak M',\allowbreak M'',\allowbreak M'''$ etc. versus limitem finitum continuo convergent, quod ita demonstramus. Supponamus primo, terminos seriei post certum intervallum continuo crescere, sive e differentiis $B-b,C-c$ etc. primam quae non evanescat
\SourceLeaf{150}
esse positivam. Sit $h$ integer talis, ut $h+b-B$ fiat quantitas positiva, statuamusque
\[
M\left(\frac{m}{m-1}\right)^h=N,\quad
M'\left(\frac{m+1}{m}\right)^h=N',\quad
M''\left(\frac{m+2}{m+1}\right)^h=N''\quad\mathrm{etc.}
\]
atque $(tt-1)^hP=Q,\ t^{2h}p=q$, ita ut $\frac{N'}{N},\frac{N''}{N'}$ etc. sint valores fractionis $\frac{Q}{q}$ ponendo $t=m,\ t=m+1$ etc. Quum itaque habeatur
\[
\begin{array}{l}
Q=t^{\lambda+2h}+At^{\lambda+2h-1}+(B-h)t^{\lambda+2h-2}\quad\mathrm{etc.}\\
q=t^{\lambda+2h}+At^{\lambda+2h-1}+bt^{\lambda+2h-2}\quad\mathrm{etc.}
\end{array}
\]
atque per hyp. $B-h-b$ sit quantitas negativa, termini seriei $N,N',N'',N'''$ etc. post certum saltem intervallum continuo decrescent, adeoque termini seriei $M,M',M'',M'''$ etc., qui illis resp. semper sunt minores, dum continuo crescunt, tantummodo versus limitem finitum convergere possunt. Q.~E.~P.

Si termini seriei $M,M',M'',M'''$ etc. post certum intervallum continuo decrescunt, accipere oportet pro $h$ integrum talem, ut $h+B-b$ sit quantitas positiva, evinceturque per ratiocinia prorsus similia, terminos seriei
\[
M\left(\frac{m-1}{m}\right)^h,\quad
M'\left(\frac{m}{m+1}\right)^h,\quad
M''\left(\frac{m+1}{m+2}\right)^h\quad\mathrm{etc.}
\]
post certum intervallum continuo crescere, unde termini seriei $M,M',M''$ etc., qui illis resp. semper sunt maiores, dum continuo decrescunt, necessario tantummodo versus limitem finitum decrescere possunt. Q.~E.~S.

IV. Denique quod attinet ad \textit{summam} seriei, cuius termini sunt $M,M',$ $M'',M'''$ etc., in casu eo, ubi hi in infinitum decrescunt, supponamus primo, $A-a$ cadere inter $0$ et $-1$, sive $A+1-a$ esse vel quantitatem positivam vel $=0$. Sit $h$ integer positivus, arbitrarius in casu eo, ubi $A+1-a$ est quantitas positiva, vel talis qui reddat quantitatem $h+m+A+B-b$ positivam in casu eo ubi $A+1-a=0$. Tunc erit
\[
\begin{array}{l}
P(t-(m+h-1))=t^{\lambda+1}+(A+1-m-h)t^\lambda\\\qquad{}+(B-A(m+h-1))t^{\lambda-1}\quad\mathrm{etc.}\\
p(t-(m+h))=t^{\lambda+1}+(a-m-h)t^\lambda\\\qquad{}+(b-a(m+h))t^{\lambda-1}\quad\mathrm{etc.}
\end{array}
\]
ubi vel $A+1-m-h-(a-m-h)$ erit quantitas positiva, vel, si haec fit $=0$, saltem $B-A(m+h-1)-(b-a(m+h))$ positiva erit. Hinc (per \textsc{i}) pro quantitate $t$ assignari poterit valor aliquis $l$, quem simulac transgressa est, valores fractionis $\frac{P(t-(m+h-1))}{p(t-(m+h))}$ semper fient positivi atque unitate maiores. Sit $n$ integer
\SourceLeaf{151}
maior quam $l$ simulque maior quam $h$, sintque termini seriei $M,M',M'',M'''$ etc. qui respondent valoribus $t=m+n,\ t=m+n+1,\ t=m+n+2$ etc., hi $N,N',N'',N'''$ etc. Erunt itaque
\[
\frac{(n+1-h)N'}{(n-h)N},\quad
\frac{(n+2-h)N''}{(n+1-h)N'},\quad
\frac{(n+3-h)N'''}{(n+2-h)N''}\quad\mathrm{etc.}
\]
quantitates positivae unitate maiores, unde
\[
N'>\frac{(n-h)N}{n+1-h},\quad
N''>\frac{(n-h)N}{n+2-h},\quad
N'''>\frac{(n-h)N}{n+3-h}\quad\mathrm{etc.}
\]
adeoque summa seriei $N+N'+N''+N'''+$ etc. maior summa seriei
\[
(n-h)N\left(\frac{1}{n-h}+\frac{1}{n+1-h}+\frac{1}{n+2-h}+\frac{1}{n+3-h}+\mathrm{etc.}\right)
\]
quotcunque termini colligantur. Sed posterior series, crescente terminorum numero in infinitum, omnes limites egreditur, quum summa seriei $1+\frac12+\frac13+\frac14+$ etc. quam infinitam esse constat etiam infinita maneat, si ab initio termini $1+\frac12+\frac13+\mathrm{etc.}$ $+\frac{1}{n-1-h}$ rescindantur. Quare summa seriei $N+N'+N''+N'''+$ etc., adeoque etiam summa huius $M+M'+M''+M'''+$ etc., cuius pars est illa, ultra omnes limites crescit.

V. Quoties autem $A-a$ est quantitas negativa absolute maior quam unitas, summa seriei $M+M'+M''+M'''+$ etc. in infinitum continuatae certo erit finita. Sit enim $h$ quantitas positiva minor quam $a-A-1$, demonstrabiturque per ratiocinia similia, assignari posse valorem aliquem $l$ quantitatis $t$, ultra quem fractio $\frac{Pt}{p(t-h-1)}$ semper adipiscatur valores positivos unitate minores. Quodsi iam pro $n$ accipitur numerus integer ipsis $l,m,h+1$ maior, terminique seriei $M,M',M'',M'''$ etc., valoribus $t=n,\ t=n+1,\ t=n+2$ etc. respondentes, designantur per $N,N',N''$ etc., erit
\[
N'<\frac{n-h-1}{n}\cdot N,\quad
N''<\frac{(n-h-1)(n-h)}{n(n+1)}\cdot N'\quad\mathrm{etc.}
\]
adeoque summa seriei $N+N'+N''+$ etc., quotcunque termini colligantur, minor producto ex $N$ in summam totidem terminorum seriei
\[
1+\frac{n-h-1}{n}+\frac{(n-h-1)(n-h)}{n(n+1)}+
\frac{(n-h-1)(n-h)(n-h+1)}{n(n+1)(n+2)}\quad\mathrm{etc.}
\]
Huius vero summa pro quolibet terminorum numero facile assignari potest; est scilicet
\[
\begin{array}{lr}
\mathrm{terminus\ primus}&=\displaystyle\frac{n-1}{h}-\frac{n-h-1}{h}
\end{array}
\]
